\section{Interoperabilidade e Transposição Clínica}

A verdadeira vocação de um ecossistema \textit{open-source} na área da saúde fracassa se este se isolar da infraestrutura hospitalar e clínica remanescente. Gêmeos digitais odontológicos operam em confluência profunda com a Internet das Coisas Médicas (IoMT), dependendo de um arcabouço rigoroso de interoperabilidade de dados e conectividade padronizada \cite{digital2025convergence}. 

Esta fluidez transacional é articulada nos seguintes padrões indispensáveis:

\begin{itemize}
    \item \textbf{Imageamento Radiográfico (DICOM):} O padrão fundamental de comunicação (Digital Imaging and Communications in Medicine) governa os maciços fluxos de tomografias CBCT. O suporte nativo pelo ITK e pelo 3D Slicer possibilita que os metadados de aquisição -- quilovoltagem, posicionamento isocêntrico, espaçamento de voxels -- sejam inteiramente lidos pelo \textit{pipeline} do gêmeo digital, preservando a coerência anatômica inicial e assegurando ausência de conversões arbitrárias deletérias.

    \item \textbf{Topologia de Superfície (STL, PLY, OBJ):} Escaneamentos intraorais (IOS) e malhas derivadas exigem manipulação precisa. Enquanto o formato STL permanece ubíquo -- mas restrito apenas à topografia geométrica em triângulos brutos -- o formato PLY incorpora dados cruciais como coloração (RGB) e opacidade das malhas. As rotinas do MeshLab e FEniCS digerem diretamente estas topologias complexas. Mais ainda, estas bibliotecas garantem que, de uma interface de planejamento, os dados sejam reencaminhados à manufatura aditiva sem corromper a fidelidade milimétrica do preparo protético.

    \item \textbf{Prontuários Eletrônicos em Saúde (EHR/PEP):} O dente não é um sistema isolado, e os reflexos sistêmicos (como periodontite induzida por diabetes tipo 2) são determinantes. A interoperabilidade estrutural com repositórios de base nacional (como o e-SUS Odonto e o padrão HL7 FHIR) é o Santo Graal do prontuário digital \cite{frota2025sus}. O pipeline do OpenCode permite assimilar perfis demográficos, anotações XML/JSON e dados de anamnese longitudinal. Quando o paciente relatar falhas metabólicas, o gêmeo reavalia a resiliência biológica em torno da ancoragem ortodôntica.

    \item \textbf{Redes Distribuídas e Tele-odontologia:} Com a regulamentação iminente da saúde digital no Brasil e resoluções do CFO sobre telemonitoramento, surge o desafio da telemetria odontológica em ambientes rurais ou descentralizados (\textit{desertos digitais}). Ecossistemas como o OpenTwins demonstraram a validade de nós federados de IoT na nuvem \cite{infante2025distributed}. Aliado ao OpenCode (com MCPs remotos \textit{stateless}), torna-se tangível simular a adaptação de tratamentos endodônticos \cite{elsaid2025digital} em clínicas periféricas do SUS, democratizando a odontologia de alto nível onde as barreiras geográficas antes a tornavam inacessível \cite{cuadros2023assessing}.
\end{itemize}
